
% ============= Packages =============

% Dokumentinformationen
\usepackage[hyphens]{url}
\usepackage[
	pdftitle={\Titel},
	pdfauthor={\Name},
	pdfkeywords={},	
	%Links nicht einrahmen
	hidelinks
]{hyperref}
% Metadaten der PDF konfigurieren

% Standard Packages
%
\renewcommand*\familydefault{\sfdefault} 
%\usepackage{lmodern}
\usepackage{helvet}
\usepackage[T1]{fontenc}
\usepackage[utf8]{inputenc}
\usepackage[english]{babel}


\usepackage{graphicx}
\graphicspath{{abbildung/}}
\usepackage[headsepline]{scrlayer-scrpage}
\usepackage{scrlayer-scrpage}
\usepackage{float}
\usepackage{relsize}

\usepackage{subcaption}
\usepackage[dvipsnames]{xcolor}
\usepackage{float}
%\usepackage{todonotes} % Geile Notizen am Rand!
% zusätzliche Schriftzeichen der American Mathematical Society
\usepackage{amsfonts}
\usepackage{amsmath,xfrac}
\setlength\jot{20pt} % Abstand der Formeln (align) innerhalb eines Absatzes 
\usepackage{diagbox} % Diagonalen Strich in Tabellen
\numberwithin{equation}{section} % Nummerierung der Gleichungen mit Kapitel.Gleichung
\sloppy
\definecolor{lightgray}{gray}{0.5}


\automark[subsection]{section}
\ihead{\Name, \Matrikelnummer}
\ohead{\leftmark}
\chead{}

% Einbinden von Matlabfiles
% WICHTIG: Die Datei mcode.sty muss sich dort befinden, wo das Hauptdokument gespeichert ist. In diesem Fall in dem Ordner wo sich die main.tex befindet.
\usepackage{listings}    
\lstset{language=Matlab,    % Einstellungen für Listing
    breaklines=true,%
    morekeywords={matlab2tikz},
    keywordstyle=\color{blue},%
    morekeywords=[2]{1}, keywordstyle=[2]{\color{black}},
    identifierstyle=\color{black},%
    stringstyle=\color{violet},
    commentstyle=\color{OliveGreen},%
    showstringspaces=false,%without this there will be a symbol in the places where there is a space
    numbers=left, %with number=none the numbers are removed
    numberstyle={\tiny \color{black}},% size of the numbers
    numbersep=9pt, % this defines how far the numbers are from the text
}
\usepackage[autolinebreaks,useliterate]{mcode}
\renewcommand{\lstlistingname}{Programmbeispiel}% Listing -> Algorithm
\renewcommand{\lstlistlistingname}{List of \lstlistingname s}

\usepackage{lscape} % Seiten im Querformat darstellen
\usepackage{qrcode} % Darstellung von QR Codes


% \usepackage{enumitem}
% \setlist{itemsep=0.2em, topsep=0.2em}

% Einbinden von PDF Dokumenten als ganze Seite
\usepackage{pdfpages}
% Tabellen die über mehr als eine Seite gehen (Formelverzeichnis)
\usepackage{longtable}

% Zeilenabstand
\usepackage{setspace}
\onehalfspacing % 1,5 facher Zeilenabstand


% Package Einstellungen & Sonstiges
% Besondere Trennungen
\hyphenation{De-zi-mal-tren-nung}

% Mathematik Eigene Befehle
\newcommand{\dt}{ \, \mathrm{d}t} % Macht ein mathematisch schönes dt

% Figure-Nr. nach Kapitel
\makeatletter
\@addtoreset{figure}{section}
\@addtoreset{table}{section}
\makeatother
\renewcommand{\thefigure}{\thesection.\arabic{figure}}
\renewcommand{\thetable}{\thesection.\arabic{table}}
